\section{The rank-one elliptic group and its local parameters}
\label{ell:section}

We prove Propositions~\ref{ell:rank} and~\ref{ell:formal}, the two elliptic
inputs used in Section~\ref{red:local-comparison-section}.  Recall the fixed
data of \eqref{ell:data}:
\[
 F=\Q(\sqrt2),\qquad d=75-53\sqrt2,\qquad
 E:\ y^2=x^3-d^2x,
\]
with \(\sigma\) the nontrivial automorphism of \(F/\Q\).  The algebraic integer \(d\)
is positive at both real embeddings and has norm \(7\).  In particular,
\(\mathfrak d=(d)\) is a prime ideal with residue field \(\F_7\).

In Section~\ref{pri:section}, we will choose split primes of good reduction
for $E$ and its conjugate $E^\sigma$ at which $7$ is a nonsquare. The identity
$dd^\sigma=7$ then makes their reductions quadratic twists with opposite
Frobenius traces.

\begin{proof}[Proof of Proposition~\ref{ell:rank}]
The ring of integers is $\mathcal O_F=\Z[\sqrt2]$.  Minkowski's bound is
$\sqrt8/2=\sqrt2<2$, so every ideal class contains an ideal of norm one.
Thus $\mathcal O_F$ is a principal ideal domain.  The unit theorem gives
$|\mathcal O_F^*/\mathcal O_F^{*2}|=4$.  The units $-1$ and $1+\sqrt2$
realize all four sign patterns at the real embeddings, so a totally
positive unit is a square.

Consider the two-isogenous curve
\[
                         E':\ y^2=x^3+4d^2x.
\]
We use the usual two-isogeny descent, with isogenies
$\varphi:E\to E'$ and $\widehat\varphi:E'\to E$ whose composite is $[2]$;
see \cite[Chapter~X, Proposition~4.7 and Example~4.9, pp.~336--337]{Silverman2009}.  Its homomorphisms
\[
 \alpha:E(F)\longrightarrow F^*/F^{*2},\qquad
 \alpha':E'(F)\longrightarrow F^*/F^{*2}
\]
send a point with nonzero $x$ to the class of $x$, send the origin to $1$,
and send $(0,0)$ to the class of the linear coefficient in the equation.
Their kernels are $\widehat\varphi E'(F)$ and $\varphi E(F)$, respectively.
In the present case
\begin{equation}\label{ell:isogeny-index}
             |\operatorname{im}\alpha|\,
             |\operatorname{im}\alpha'|=4\,2^{\rank E(F)}.
\end{equation}
For completeness, $E$ has full rational two-torsion, and either point
$(\pm d,0)$ maps under $\varphi$ to the nonidentity point of
$\ker\widehat\varphi$.  Hence the kernel of
$E'(F)/\varphi E(F)\to \widehat\varphi E'(F)/2E(F)$ is trivial.  Taking
indices in $2E(F)\subset\widehat\varphi E'(F)\subset E(F)$ gives the left
side of \eqref{ell:isogeny-index} as $|E(F)/2E(F)|$, which is
$2^{\rank E(F)+2}$ by the Mordell--Weil Theorem.

If the linear coefficient of an equation $y^2=x^3+cx$ is a unit at a
finite place, then a nonzero $x$ has even valuation there.  Indeed, when
$v(x)<0$ the term $x^3$ has strictly smaller valuation than $cx$, while
when $v(x)>0$ the term $cx$ has strictly smaller valuation; an odd
$v(x)$ would give an odd valuation of $y^2$ in either case.  Applied to
$E$, this shows that every class in $\operatorname{im}\alpha$ is supported,
apart from its unit class, only at $\mathfrak d$.  The class number
calculation therefore bounds the image by $4\cdot2=8$.  The torsion
points already give the four distinct classes
\[
                             1,-1,d,-d.
\]
Let $w=-5+4\sqrt2=d(1+\sqrt2)^3$.  Direct calculation gives
$w^3-w=-8d$, and consequently
\[
                  Q=(-dw,\,2\sqrt2\,d^2)\in E(F).
\]
The element $-dw$ has opposite signs at the two real embeddings, whereas
the four torsion classes have equal signs at both embeddings.  Thus
$\alpha(Q)$ is outside those four classes, proving
$|\operatorname{im}\alpha|=8$.

We next prove $\operatorname{im}\alpha'=\{1\}$.  At both real embeddings,
$x^3+4d^2x$ has the sign of $x$, so every nonzero $x$ on $E'$ is totally
positive.  The unit-coefficient argument excludes odd valuations away
from the prime above $2$ and $\mathfrak d$.  Normalize the valuation at
the prime above $2$ to have value group $\Z$.  Then $v(4d^2)=4$.  If
$v(x)$ were odd, the valuations $3v(x)$ and $4+v(x)$ would be unequal and
their minimum would be odd, again impossible for $y^2$.
At $\mathfrak d$ the coefficient has valuation $2$.  An odd $v(x)$ is
impossible unless $v(x)=1$, by the same comparison.  In that remaining
case write $x=du$ with $u$ a local unit.  The equation becomes
\[
                         y^2=d^3u(u^2+4).
\]
An even valuation would require $u^2+4\equiv0\pmod{\mathfrak d}$.
But $-4=3$ is not a square in $\F_7$.  Thus all finite valuations of $x$
are even.  Since $\mathcal O_F$ is principal, $x$ is a square times a
unit; total positivity makes that unit a square.  The exceptional point
$(0,0)$ also gives the trivial class $4d^2$, and the origin does likewise.
This proves the assertion about $\alpha'$.  Equation
\eqref{ell:isogeny-index} now gives $\rank E(F)=1$.

Finally, choose $m$ divisible by the exponent of the finite torsion
subgroup of $E(F)$.  The finitely generated group $mE(F)$ is torsion-free
of rank one and therefore has a generator $P$ as claimed.
\end{proof}

Recall the parameter \(z=-x/y\), with \(z(O)=0\), and the formal
logarithm and its degree-\(K\) truncation
\[
 \ell(Z)=Z+\sum_{t\ge2}c_tZ^t,\qquad
 \ell_K(Z)=Z+\sum_{2\le t\le K}c_tZ^t.
\]

Reading multiplication indices from elliptic coordinates has a
precedent in Poonen's rank-one Diophantine construction.  His
Lemma~9 reads divisibility of indices from denominators of elliptic
multiples, and Lemma~11 uses the formal group to approximate an
integer square by a quotient of coordinates
\cite{PoonenRankOne2002}.  Proposition~\ref{ell:formal} uses a
truncated formal logarithm to approximate a quotient of indices.
The paired local tests in the reduction use this comparison together with the
conjugate curve.

The effective computation of the truncations in Section~\ref{red:local-comparison-section} can be made explicit.  With $w=-1/y$ one has
\[
              x=z/w,\qquad w=z^3-d^2zw^2.
\]
Successively solving this identity determines $w$ as a formal power
series in $z$; substitution in $dx/(2y)$ and termwise integration
determines $\ell$.  The formal group and its invariant differential have
integral coefficients at every place of good integral reduction
\cite[Chapter IV]{Silverman2009}.

\begin{proof}[Proof of Proposition~\ref{ell:formal}]
Include $2$, the primes of bad reduction of either equation, and the
ramified primes of $F/\Q$ in $S_E$.  Let $p\notin S_E$, let $v$ be a
place above $p$ normalized by $v(p)=1$, and let $n,n'\in\Z$, with
$n'\ne0$ and both $nP,n'P$ reducing to $O$, as in the proposition.
At this allowed place $F_v/\Q_p$ is
unramified, so its normalized value group is $\Z$.  The kernel of
reduction is the formal group on the maximal ideal.  For a finite point
on this integral good model, membership in that kernel is equivalent to
$v(x)<0$; its parameter then has $v(z)>0$.

Write the invariant differential as
$(1+\sum_{j\ge1}a_jZ^j)dZ$, with $a_j$ integral at $v$.  Thus
$v(c_t)\ge-v_p(t)$.  If $v(Z)\ge1$, the term of degree $t\ge2$ in the
logarithm has valuation at least $t\,v(Z)-v_p(t)>v(Z)$ for $p>2$, and
these valuations tend to infinity.  The logarithm consequently converges
on the entire formal kernel, is a homomorphism there, and satisfies
\[
                         v(\ell(Z))=v(Z).
\]
These are also the standard formal-logarithm assertions of
\cite[Chapter IV, \S6]{Silverman2009}; the displayed coefficient bound
specifies the uniformity we require.

Set $U=nP$ and $V=n'P$.  The identity $[n']U=[n]V$, entirely within
the formal kernel, gives $n'\ell(z(U))=n\ell(z(V))$.
The point $V$ is nonzero, hence $\ell(z(V))\ne0$.  Valuation preservation
proves \eqref{ell:ratio}, including $U=O$.

Now fix $K\ge1$ and suppose $p>K+1$ and $v_p(n/n')\ge0$.  For every $t>K$,
\begin{equation}\label{ell:tail-exponent}
                           t-1-v_p(t)\ge K.
\end{equation}
For $t<p$ this is immediate.  For $t\ge p$, if $j=v_p(t)>0$ then
$t-1-j\ge p^j-1-j\ge p-2\ge K$; if $j=0$ it is again immediate.
Therefore, for $v(Z)\ge1$,
\begin{equation}\label{ell:tail}
                 v(\ell(Z)-\ell_K(Z))\ge v(Z)+K.
\end{equation}
The inequality also holds at $Z=0$ with the stated convention.  Put
$q=n/n'$.  The hypothesis $v_p(q)\ge0$ and \eqref{ell:ratio} give
$v(z(U))\ge v(z(V))$.  The identity $\ell(z(U))=q\ell(z(V))$ and
\eqref{ell:tail} imply
\[
          v\bigl(\ell_K(z(U))-q\ell_K(z(V))\bigr)
                                    \ge v(z(V))+K.
\]
The same tail estimate gives
$v(\ell_K(z(V)))=v(z(V))$, in particular a nonzero denominator.
Division proves \eqref{ell:truncation}.  The proof applies verbatim to
the conjugate equation and the polynomial $\ell_K^\sigma$.
\end{proof}

The quantifiers in Proposition~\ref{ell:formal} are those used in
Lemma~\ref{red:local-comparison}: for each fixed \(K\), a single finite
set of excluded rational primes works for all integer indices and all
unramified places above the remaining primes.  The transfer there concerns
integer multiples of the fixed \(F\)-rational points \(P,P^\sigma\)
and finite polynomial evaluations.
No infinite series is evaluated in the elementary extension.
